(1) 对 $m=0$，平均透射幅度为 $c_0=(1+e^{i\phi})/2=e^{i\phi/2}\cos(\phi/2)$。对 $m\ne0$，两半周期积分相减给 $c_m=(e^{i\phi}-1)(1-(-1)^m)/(i2\pi m)$。因此非零偶数级消失，而奇数 $m$ 的振幅为 $c_m=2e^{i\phi/2}\sin(\phi/2)/(\pi m)$。

(2) 奇数级的强度比例为 $|c_m|^2=4\sin^2(\phi/2)/(\pi^2m^2)$。利用所有正负奇整数的平方倒数和 $\sum_{m\ne0,\ m\ \rm odd}m^{-2}=\pi^2/4$，所以 $\sum_{m\ne0}|c_m|^2=\sin^2(\phi/2)$；加上零级 $|c_0|^2=\cos^2(\phi/2)$ 恰为1，体现无损单位模透射。

(3) 第一对级次 $m=\pm1$ 的总效率 $\eta_{\pm1}=2\times4\sin^2(\phi/2)/\pi^2=8\sin^2(\phi/2)/\pi^2$。在 $\phi=(2j+1)\pi$ 时取最大 $8/\pi^2\approx81.1\%$，此时零级完全压制，剩余能量进入更高奇数级。

(4) 透射光栅给横向波矢 $k_{x,m}=2\pi m/d$，所以 $\sin\theta_m=m\lambda/d$。仅当 $|m|\lambda\le d$ 时相应级次成为远场传播光。傅里叶系数归一用于理想薄屏近轴标量传播；高数值孔径下，斜向能流权重和倏逝场需通过严格Maxwell衍射重算。
